% ======================================================================
% sections/abstract.tex
%
% Abstract scaffold for the paper:
%   Triple Humanoid 24/7 Adverse Event Oncology Trial Response Team:
%   4-Site Rotation
%
% This file is a draft scaffold. A future Claude Code Opus 4.7 1M Max
% session reads the bracketed instruction blocks below, then reads the
% files those instructions name, then expands the scaffold into the
% finished abstract text. The future session should overwrite both the
% bracketed instructions and the placeholder prose with the finished
% paragraph, while keeping the abstract length to roughly 300 words.
% ======================================================================

\begin{sectionaccent}

[ABSTRACT INSTRUCTIONS FOR FUTURE CLAUDE CODE OPUS 4.7 1M MAX SESSION.
The abstract has two halves and reads as a single continuous paragraph
of about 300 words. The first half is the production overview: it must
walk the reader through the full chain of code instructions, code
generations, code execution, draft paper generation, and full paper
production. The second half is the results catch: it must surface
specific numerical results, discussion claims, and metrics that appear
elsewhere in the paper so that a reader who only reads the abstract
still walks away with concrete numbers. Do not over generalize. Use
the file list below.

Files to read before writing the production overview half:
\begin{itemize}
  \item \path{demo-projects/07-humanoid/paper/instructions/README.md}
        for the project thesis, the four site identifiers (SF-01, SD-01,
        BO-01, AT-01), the 12 H2 robot total count, the 1 Hz LLM
        broadcast cadence, the 10 Hz humanoid motion cadence, the 168
        hour monitoring window, and the camaraderie reduction factor of
        3 in single robot error potential.
  \item \path{demo-projects/07-humanoid/paper/instructions/00-project-overview/}
        through \path{09-runtime-instructions/} for the 10 instruction
        directories that drive code generation.
  \item \path{demo-projects/07-humanoid/paper/codegen/README.md} and
        \path{demo-projects/07-humanoid/paper/codegen/architecture.md}
        for the code generation phase output.
  \item \path{demo-projects/07-humanoid/paper/execution/README.md}
        and \path{demo-projects/07-humanoid/paper/execution/reports/execution_report_v0_5_0.md}
        for the execution phase output.
  \item \path{demo-projects/07-humanoid/paper/inputs/README.md},
        \path{demo-projects/07-humanoid/paper/inputs/Deep-Research-A/README.md},
        and the equivalent README in
        \path{demo-projects/07-humanoid/paper/inputs/Deep-Research-B/}
        for the prior paper context (the prior paper covered Medical
        Full Automation A and Medical Rapid Prototyping B).
\end{itemize}

Files to read before writing the results catch half:
\begin{itemize}
  \item \path{demo-projects/07-humanoid/paper/codegen/data/iterations/index.jsonl}
        for the 32 deterministic iteration sweep results across the 4
        sites and many columns. Quote the iteration count and the
        per-site count in the abstract.
  \item \path{demo-projects/07-humanoid/paper/codegen/reports/dashboard.html}
        and \path{demo-projects/07-humanoid/paper/codegen/reports/report.md}
        for the comparison summary headline numbers. Surface the lead
        metric (the camaraderie factor of 3, plus the e-stop latency of
        5 ms swarm wide).
  \item \path{demo-projects/07-humanoid/paper/execution/data/site_runtime_runs/AT-01_llm_decisions.jsonl}
        plus \texttt{BO-01}, \texttt{SD-01}, and \texttt{SF-01} variants
        in the same directory for the per-site LLM decision streams at
        1 Hz. Note that the 4 site files are identical in byte length
        (62,330 bytes each); call this out as a check that the runtime
        produced parallel-scope decision logs across sites and not as a
        data duplication artifact in production.
  \item \path{demo-projects/07-humanoid/paper/execution/data/iterations_aggregate.duckdb}
        for the DuckDB roll up of all 32 iteration sweeps.
\end{itemize}

Required content cues for the abstract paragraph:
\begin{itemize}
  \item Open with the thesis sentence verbatim from the paper README.
  \item Name the 4 sites (SF-01, SD-01, BO-01, AT-01), the 3 H2 per
        site, the 12 total H2, and the 1 on-prem Claude Opus 4.7 1M
        instance per site.
  \item Name the 1 Hz broadcast cadence and the 10 Hz robot motion
        cadence.
  \item Name the 168 hour monitoring window and the approximate 84 AE
        plus 24 SAE event volume.
  \item Name CTCAE version 5 and the FDA RTCT 1 hour SLA from HR 9505
        as the regulatory anchor.
  \item Name the 32 iteration sweeps and the camaraderie reduction
        factor of 3 in single robot error potential.
  \item Cite the author's December 2025 paper
        (\path{kawchak_2025_18029100}) as the prior evidence that
        LLMs are effective on FDA adverse events data.
  \item Close on the future industry value: the on-prem repository LLM
        plus sensor pipeline plus Cartesian x, y, z command structure
        is presented as a checklist already executed so the industry
        does not have to repeat it.
\end{itemize}

Do not exceed roughly 300 words. Do not introduce numbers that are not
supported by the files named above. Do not use em dashes, double
dashes, or triple dashes. Replace any stray \texttt{SS} with the
section sign \S\ during the cleanup pass.]

The abstract will summarize the on-premises repository based LLM
pipeline, the 4-site rotation across SF-01, SD-01, BO-01, and AT-01,
the 3 Unitree H2 EDU robots per site, the 1 Hz broadcast cadence, the
Cartesian x, y, z command surface, and the per-site camaraderie
behavior that reduces single robot error potential. The closing
sentences will surface specific results from the 32 iteration sweeps
and the per-site LLM decision streams.

\sectiontable{Abstract planning matrix}{%
Thesis & A1 & 1 & instructions README & State thesis verbatim & 1 sentence & Reads like the README & Ready \\
Overview & A2 & 2 & code, exec README & Production chain & 4 to 5 sent. & Names phases & Ready \\
Numbers & A3 & 3 & index, dashboard & Headline metrics & 3 to 4 sent. & Sourced numbers & Ready \\
}

\end{sectionaccent}
